\documentclass[UTF8]{ctexart}
\usepackage{amsmath}
\usepackage{booktabs}
\usepackage{geometry}
\usepackage{listings}
\usepackage{xcolor}
\usepackage{hyperref}
\usepackage{graphicx}

\geometry{a4paper, left=2.5cm, right=2.5cm, top=2.5cm, bottom=2.5cm}
\hypersetup{colorlinks=true, linkcolor=blue, urlcolor=blue, citecolor=blue}

\lstset{
  basicstyle=\ttfamily\small,
  breaklines=true,
  frame=single,
  keywordstyle=\color{blue},
  commentstyle=\color{gray},
  stringstyle=\color{teal},
  showstringspaces=false,
  tabsize=2,
  columns=flexible,
}

\title{\textbf{系统工具开发课程实验报告（第一周）}}
\author{廖世嘉 \quad 25020007073 \\ 25级计算机科学与技术}
\date{2026年8月24日}

\begin{document}
\maketitle
\tableofcontents
\newpage

\section{实验基本信息}

\begin{table}[htbp]
\centering
\caption{实验基本信息}
\label{tab:info}
\begin{tabular}{ll}
\toprule
项目 & 内容 \\
\midrule
姓名 & 廖世嘉 \\
学号 & 25020007073 \\
专业 & 25级计算机科学与技术 \\
实验环境 & Linux / Bash / Git / XeLaTeX \\
实验日期 & 2026年8月24日 \\
GitHub仓库 & \href{https://github.com/MrLLLiao/SystemToolkitDevCourse}{https://github.com/MrLLLiao/SystemToolkitDevCourse} \\
\bottomrule
\end{tabular}
\end{table}

\section{实验目的}

本实验围绕 Shell 基础与文件系统、Shell 管道与文本处理、Git 版本控制以及 LaTeX 文档编辑四个主题展开。通过四个连续实验，熟悉命令行环境下的文件操作、脚本编写、日志分析、分支管理、冲突解决和技术文档构建方法。同时，通过命令行验证文件权限、脚本退出码、Git 工作区状态以及 PDF 交叉引用，形成完整的"操作——验证——排错"实践流程。

\section{实验环境与工具}

本实验在 Linux 环境中完成，主要使用 Bash、find、ls、cp、chmod、awk、sort、head、Git、XeLaTeX/latexmk、pdftotext 与 file 等工具。实验目录为 SystemToolkitDevCourse，四个实验分别位于 q01、q02、q03 和 q04。

\section{实验一：含空格文件名的批量整理}

\subsection{实验要求}

在 q01 中创建 input/docs 和 input/tmp，并创建 notes one.txt（内容为 alpha 和 beta 两行）、隐藏文件 .secret.txt（内容为 hidden）、空文件 empty.txt 以及两行日志。随后将所有 .txt 文件复制到 work/25020007073，保持相对目录结构，并将目录权限设为 750、普通文件权限设为 640，最后生成 inventory.txt 列出相对路径和字节数。

\subsection{主要操作}

使用 mkdir -p 创建目录，利用 printf 和 touch 创建文件。对于包含空格的文件名，在 Shell 中使用引号包围整个路径：

\begin{lstlisting}[language=bash]
printf "alpha\nbeta\n" > "input/docs/notes one.txt"
printf "hidden\n" > input/docs/.secret.txt
touch input/tmp/empty.txt
printf "INFO: program started\nINFO: program finished\n" > input/run.log
\end{lstlisting}

使用 ls -laR input 以长格式递归显示全部项目，其中 -a 保证隐藏文件能够显示。复制文件时使用 find ... -print0 配合 read -d ''，避免文件名中的空格导致错误分割：

\begin{lstlisting}[language=bash]
STUDENT_ID="25020007073"
mkdir -p "work/$STUDENT_ID"
find input -type f -name "*.txt" -print0 |
while IFS= read -r -d '' file; do
  relative="${file#input/}"
  mkdir -p "work/$STUDENT_ID/$(dirname "$relative")"
  cp -- "$file" "work/$STUDENT_ID/$relative"
done
\end{lstlisting}

然后使用 find 配合 chmod 设置权限：

\begin{lstlisting}[language=bash]
find "work/$STUDENT_ID" -type d -exec chmod 750 {} +
find "work/$STUDENT_ID" -type f -exec chmod 640 {} +
\end{lstlisting}

生成清单：

\begin{lstlisting}[language=bash]
find "work/$STUDENT_ID" -type f ! -name "inventory.txt" \
  -printf '%P\t%s\n' > "work/$STUDENT_ID/inventory.txt"
chmod 640 "work/$STUDENT_ID/inventory.txt"
\end{lstlisting}

\subsection{结果与验证}

实际检查结果表明，input 下正确存在 notes one.txt、.secret.txt、empty.txt 和 run.log。其中前三个文本文件的大小分别为 11、7 和 0 字节。复制后的 work/25020007073 中，目录显示为 drwxr-x---，即 750；普通文件显示为 -rw-r-----，即 640。inventory.txt 成功记录复制后的相对路径和字节数：

\begin{lstlisting}
tmp/empty.txt 0
docs/.secret.txt 7
docs/notes one.txt 11
\end{lstlisting}

\section{实验二：随机访问日志统计}

\subsection{实验要求}

在 q02 中创建给定的 CSV 数据并编写 analyze.sh。脚本接收 CSV 路径，若文件不存在则将错误写入标准错误并返回非零退出码；正常情况下使用 Shell 管道、awk、sort、head 统计 HTTP 5xx 次数最多的前两个 path，并计算全部数据行的平均延迟。

\subsection{脚本实现}

文件检查部分如下：

\begin{lstlisting}[language=bash]
file="$1"
if [ -z "$file" ]; then
  echo "Error: missing CSV file path" >&2
  exit 1
fi
if [ ! -f "$file" ]; then
  echo "Error: file not found: $file" >&2
  exit 1
fi
\end{lstlisting}

5xx 统计通过 awk 对状态码进行过滤，再通过 sort 和 head 完成排序与截取：

\begin{lstlisting}[language=bash]
awk -F',' 'NR > 1 && $4 ~ /^5[0-9][0-9]$/ { count[$3]++ }
END { for (path in count) print count[path], path }' "$file" |
sort -k1,1nr -k2,2 |
head -n 2
\end{lstlisting}

sort -k1,1nr 负责按次数数值降序排列，-k2,2 在次数相同时按路径字典序排列。平均延迟通过另一个 awk 忽略表头并保留两位小数：

\begin{lstlisting}[language=bash]
awk -F',' 'NR > 1 { sum += $5; count++ }
END { printf "%.2f\n", sum / count }' "$file"
\end{lstlisting}

\subsection{结果}

运行 ./analyze.sh access.csv 得到：

\begin{lstlisting}
Top 2 5xx paths:
2 /api/orders
2 /api/users
Average latency_ms:
193.75
\end{lstlisting}

运行不存在的文件时，脚本输出文件不存在错误并返回退出码 1，说明标准错误和非零退出状态处理符合要求。本实验未使用 Python、电子表格或手工统计。

\section{实验三：制造、解决并解释一次 Git 合并冲突}

\subsection{实验要求}

在 q03 中从零初始化仓库。在 main 中创建包含 mode=normal 的 config.txt 并提交；然后创建 feature-a 改为 mode=safe 并提交；从初始 main 创建 feature-b 改为 mode=fast 并提交。回到 main 后先合并 feature-a，再合并 feature-b，主动制造冲突并解决。

\subsection{实验过程}

初始化仓库：

\begin{lstlisting}[language=bash]
git init
git branch -M main
printf "mode=normal\n" > config.txt
git add config.txt
git commit -m "Initial commit"
\end{lstlisting}

创建并提交 feature-a：

\begin{lstlisting}[language=bash]
git switch -c feature-a
printf "mode=safe\n" > config.txt
git add config.txt
git commit -m "Set mode to safe"
\end{lstlisting}

回到初始 main 后创建并提交 feature-b：

\begin{lstlisting}[language=bash]
git switch main
git switch -c feature-b
printf "mode=fast\n" > config.txt
git add config.txt
git commit -m "Set mode to fast"
\end{lstlisting}

回到主分支并按要求合并：

\begin{lstlisting}[language=bash]
git switch main
git merge feature-a
git merge feature-b
\end{lstlisting}

\subsection{冲突原因与解决}

feature-a 将初始状态 mode=normal 改为 mode=safe，而 feature-b 将同一行改为 mode=fast。当两个修改合并时，Git 无法自动判断应保留哪一个版本，因此要求人工处理冲突。本实验按要求保留 mode=safe，并增加 note=reviewed：

\begin{lstlisting}
mode=safe
note=reviewed
\end{lstlisting}

随后执行：

\begin{lstlisting}[language=bash]
git add config.txt
git commit -m "Merge feature-b and resolve conflict"
\end{lstlisting}

最终使用 git status 检查，工作区显示 nothing to commit, working tree clean；使用 git log --all --graph --decorate --oneline 能够观察到初始提交、两个特性分支提交以及最终合并提交，说明版本历史完整。

\section{实验四：修复并构建一页技术说明}

\subsection{实验要求}

在 q04/report.tex 中加入带编号和标签的公式，在正文中使用 \textbackslash ref 或 \textbackslash eqref 引用；加入 2 行 2 列表格并设置 caption 和 label；使用 latexmk -pdf -halt-on-error 或 tectonic 生成 PDF，并确认交叉引用没有显示"??"。

\subsection{公式与交叉引用}

实验采用如下公式表示总延迟：

\begin{equation}
L_{\mathrm{total}} = \sum_{i=1}^{n} L_i
\label{eq:latency}
\end{equation}

正文通过式~\eqref{eq:latency} 进行引用。由于公式包含 equation 环境和 label，LaTeX 能够为公式自动编号并建立交叉引用。

\subsection{表格与交叉引用}

实验中使用严格的 2 行 2 列表格：

\begin{table}[htbp]
\centering
\caption{日志统计结果示例}
\label{tab:result}
\begin{tabular}{ll}
\toprule
Item & Value \\
\midrule
Average latency & 193.75 ms \\
\bottomrule
\end{tabular}
\end{table}

正文通过表~\ref{tab:result} 引用该表，表格具有独立的编号、标题和标签。

\subsection{构建与验证}

原实验中使用 pdflatex 编译含中文作者信息的文档时出现 Unicode 字符错误。原因是 pdfLaTeX 对中文 Unicode 字符的直接支持有限。修正后采用 CTeX/XeLaTeX 方案生成文档。使用命令行构建：

\begin{lstlisting}[language=bash]
latexmk -pdf -halt-on-error report.tex
\end{lstlisting}

生成 PDF 后使用 pdftotext 和 grep 检查交叉引用，结果显示 PDF 文件能够正常生成，正文中的公式引用显示为对应编号，表格引用显示为"Table 1"，而 grep -Fn '??' 无输出，说明交叉引用没有出现未解析的"??"。

\section{实验结果汇总}

\begin{table}[htbp]
\centering
\caption{四个实验的完成情况}
\label{tab:summary}
\begin{tabular}{cll}
\toprule
题号 & 实验主题 & 验证结果 \\
\midrule
1 & Shell 文件系统 & 空格文件名、隐藏文件、复制、权限和 inventory 均正确 \\
2 & Shell 日志分析 & 前两名 5xx path 及平均延迟统计正确，异常文件返回非零状态 \\
3 & Git 版本控制 & 成功产生、解决冲突，最终工作区干净并可查看图形历史 \\
4 & LaTeX 文档构建 & PDF 成功生成，公式和表格交叉引用正常，无"??" \\
\bottomrule
\end{tabular}
\end{table}

\section{问题与解决方法}

\subsection{Windows 与 Linux 权限命令差异}

实验初期曾在 Windows PowerShell 中执行 chmod，系统提示命令不存在。由于题目要求使用 750 和 640 这类 Unix 权限模式，因此将相关操作放到 Linux 环境下，通过 find 和 chmod 完成。

\subsection{LaTeX 中文编译问题}

使用 pdfLaTeX 直接编译包含中文姓名的文档时出现 Unicode 字符错误。之后改用 CTeX/XeLaTeX，使中文文本和标题能够正常排版。这一问题体现了不同 TeX 引擎在字体和 Unicode 支持方面的差异。

\subsection{交叉引用检查方法}

检查 PDF 文本中的"??"时，应采用固定字符串搜索而不是容易产生歧义的正则表达式，因此最终使用 grep -Fn '??' 进行验证。

\section{实验总结}

通过本次实验，我系统练习了 Linux 命令行环境中的文件管理、权限控制、Shell 管道、文本分析、Git 分支管理和 LaTeX 文档构建。实验一使我熟悉了隐藏文件、空格文件名和权限模式；实验二让我理解了多个 Unix 文本工具通过标准输入输出构成流水线的方式；实验三让我实际经历了 Git 合并冲突从产生到解决的全过程；实验四则强化了技术文档的结构化编辑、交叉引用和自动化构建能力。

本次实验中最重要的认识是：命令行任务不仅要"做出来"，还应当通过命令验证最终状态。例如，文件权限需要通过 ls -l 或 find 检查，Shell 脚本需要检查退出码，Git 需要通过 git status 和日志确认状态，LaTeX 则需要检查生成的 PDF 与交叉引用。通过这些验证步骤，可以显著降低操作错误并提高实验结果的可重复性。

\end{document}
